% Л8 -- ASP.NET Core: хост и конвейер запроса. Slides.
%
% Companion to the spoken script in ../L08-host-and-pipeline.md, which is the
% source of truth: the slides carry diagrams, code, real terminal output and
% short marks, and the delivery lives in the script only -- this deck carries
% no speaker notes. Slide text is in Russian because that is the content;
% everything else in this file is English.
%
% Build:  make          (latexmk -pdf; engine: pdflatex)
%
% Every C# listing was compiled and run on dotnet SDK 10.0 in a stock
% `dotnet new web -n HelloWorld`; terminal output on the slides is captured
% from those runs, with the home directory replaced by /home/user. See
% ../CLAUDE.md for the verification table.
%
% House rules: ../CLAUDE.md and ../../../CLAUDE.md. Theme: ../../../../theme/.

\documentclass[aspectratio=169,10pt,t]{beamer}

\usetheme{dotnet}
\usetikzlibrary{arrows.meta,positioning,calc,decorations.pathreplacing,fit}
\usepackage{array}   % ragged-right table columns

% A ragged-right paragraph column: justified cells in narrow columns hyphenate
% badly, and a split word is the one thing the back row cannot read.
\newcolumntype{L}[1]{>{\raggedright\arraybackslash}p{#1}}
\newcommand{\nohyph}{\hyphenpenalty=10000\exhyphenpenalty=10000\relax}

\dotnetlecture{Л8}

\title[Л8. Хост и конвейер]{ASP.NET Core: хост и конвейер запроса}
\author{Дмитрий Курочкин}
\date{}

% Code sizes, one step down from \footnotesize. Smaller than 6.6pt does not
% carry to the back of the room.
\newcommand{\codesm}{\ttfamily\fontsize{7.6}{9.0}\selectfont}
\newcommand{\codexs}{\ttfamily\fontsize{7.0}{8.3}\selectfont}
\newcommand{\termsm}{\ttfamily\fontsize{6.6}{7.9}\selectfont}

% A mark at the start of a listing line, for labels set beside the listing.
% Used through the listing's escapechar; the labels are drawn by an overlay
% picture after the listing, so the build needs two passes (latexmk does that).
\newcommand{\mk}[1]{\tikz[remember picture,overlay]\coordinate (#1) at (0,0.55ex);}

\tikzset{
  rawbrace/.style={decorate,decoration={brace,amplitude=4pt},hairline,
                   line width=0.8pt},
  rawlab/.style={anchor=west,font=\footnotesize,inner sep=0pt},
  box/.style={draw=hairline,line width=0.8pt,fill=codebg,inner sep=4pt,
              align=center,minimum height=0.9cm},
  arr/.style={-{Straight Barb[length=1.6mm]},muted,line width=0.7pt},
  arra/.style={-{Straight Barb[length=1.8mm]},accent,line width=0.9pt},
  arrl/.style={-{Straight Barb[length=1.8mm]},alt,line width=0.9pt},
}

\begin{document}
\shorthandoff{"}

% ------------------------------------------------------------------ title --
% 1
\begin{frame}[plain]
  \titlepagednc{Лекция 8}{ASP.NET Core: хост и конвейер запроса}{}{Дмитрий
    Курочкин}
\end{frame}

% ============================================== 1. ENTRY POINT AND LIFETIME ==
\section{1. Точка входа и жизненный цикл}

% 2
\begin{frame}[plain]
  \sectionopener{1}{Точка входа и жизненный цикл}{Шесть строк, и в них
    \hi{вся} структура сервиса.}
\end{frame}

% 3
\begin{frame}[plain]
  \vote{Сколько строк кода нужно, чтобы поднять на .NET HTTP-сервер, который
    отвечает на запрос?}
\end{frame}

% 4
\begin{frame}[fragile]{Шесть строк}
  \slidetop
  \claim{Каждая строка --- \hi{фаза} жизни сервиса.}

\begin{lstlisting}[style=csharp,basicstyle=\ttfamily\fontsize{8.4}{10.2}\selectfont,escapechar=|]
|\mk{a1}|var builder = WebApplication.CreateBuilder(args);

|\mk{a2}|builder.Services.AddHealthChecks();

|\mk{a3}|var app = builder.Build();

|\mk{a4}|app.MapHealthChecks("/health");
|\mk{a5}|app.MapGet("/", () => "Hello from BSUIR");

|\mk{a6}|app.Run();
\end{lstlisting}
  \begin{tikzpicture}[remember picture,overlay]
    \node[rawlab] at ([xshift=8.9cm]a1) {\textbf{1} создать builder};
    \node[rawlab] at ([xshift=8.9cm]a2) {\textbf{2} зарегистрировать сервисы};
    \node[rawlab,text=accent] at ([xshift=8.9cm]a3) {\textbf{3} собрать приложение};
    \draw[rawbrace] ([xshift=8.65cm,yshift=0.35em]a4) -- ([xshift=8.65cm,yshift=-0.3em]a5);
    \node[rawlab] at ([xshift=8.9cm]$(a4)!0.5!(a5)$) {\textbf{4} собрать конвейер};
    \node[rawlab] at ([xshift=8.9cm]a6) {\textbf{5} запустить};
  \end{tikzpicture}

  \slidegap
  \aside{Порядок фаз --- то, что нужно понимать. Всё остальное в сервисе
    встраивается в эти строки.}
\end{frame}

% 5
\begin{frame}{\texttt{CreateBuilder}: что уже готово}
  \slidetop
  \claim{Одна строка --- и пять вещей \hi{уже} собраны.}

  {\renewcommand{\arraystretch}{1.5}
  \begin{tabular}{@{}L{3.6cm}L{9.2cm}@{}}
    \hi{\textbf{конфигурация}} & файлы настроек, переменные окружения, аргументы
      командной строки \\
    \hi{\textbf{логирование}}  & в консоль, уровни --- из настроек \\
    \hi{\textbf{контейнер}}    & зависимостей: создан, пока пустой \\
    \hi{\textbf{окружение}}    & Development, Staging или Production \\
    \hi{\textbf{сервер}}       & Kestrel: зарегистрирован, ещё не слушает \\
  \end{tabular}}

  \slidegap
  \aside{Вы этого не писали: всё приехало с \texttt{CreateBuilder}.}
\end{frame}

% 6
\begin{frame}[fragile]{Откуда конфигурация}
  \slidetop
  \claim{Позднее \hi{перекрывает} раннее.}

  \begin{center}
    \begin{tikzpicture}[font=\footnotesize,
      cb/.style={box,minimum width=2.75cm,minimum height=1.0cm}]
      \node[cb] (f1) at (0,0)    {\texttt{appsettings.json}};
      \node[cb] (f2) at (3.25,0) {\texttt{appsettings.}\\\texttt{Development.json}};
      \node[cb] (f3) at (6.5,0)  {переменные\\окружения};
      \node[cb,draw=accent,line width=1.0pt] (f4) at (9.75,0) {\texttt{args}};
      \draw[arr] (f1) -- (f2);
      \draw[arr] (f2) -- (f3);
      \draw[arr] (f3) -- (f4);
      \node[font=\scriptsize,text=muted,below=0.25em] at (f1.south) {умолчание};
      \node[font=\scriptsize,text=accent,below=0.25em] at (f4.south) {побеждает};
    \end{tikzpicture}
  \end{center}

  \slidegap
\begin{lstlisting}[style=plain,basicstyle=\codexs]
# appsettings.json      "Logging": { "LogLevel": { "Default": "Information" } }
# environment variable  Logging__LogLevel__Default=Warning
# command line          dotnet run --Logging:LogLevel:Default=Debug
\end{lstlisting}

  \slidegap
  \aside{Окружение --- переменная \texttt{ASPNETCORE\_ENVIRONMENT}. Не задана
    --- приложение считает себя \lo{Production}: без подробных ошибок наружу.}
\end{frame}

% 7
\begin{frame}[fragile]{Регистрация сервисов}
  \slidetop
  \claim{Регистрация --- это \hi{описание}, а не создание.}

\begin{lstlisting}[style=csharp,basicstyle=\ttfamily\small]
builder.Services.AddHealthChecks();
builder.Services.AddAuthentication();
builder.Services.AddAuthorization();
\end{lstlisting}

  \slidegap
  \begin{cbul}{accent}
    \item ни один объект ещё \hi{не создан}
    \item \texttt{Add...} --- метод расширения, внутри десяток регистраций
    \item файл запуска читается как список \hi{модулей}, а не классов
  \end{cbul}
\end{frame}

% 8
\begin{frame}{\texttt{Build}}
  \slidetop
  \claim{Здесь описание \hi{становится} приложением.}

  {\renewcommand{\arraystretch}{1.55}
  \begin{tabular}{@{}L{3.6cm}L{9.2cm}@{}}
    \hi{\textbf{контейнер}}    & заморожен: зарегистрировать больше нельзя \\
    \hi{\textbf{конфигурация}} & зафиксирована \\
    \hi{\textbf{сервер}}       & создан, но ещё не слушает \\
    \hi{\textbf{app}}          & умеет две вещи: собирать конвейер и запускаться \\
  \end{tabular}}

  \slidegap
  \aside{Не ограничение, а гарантия: всё, что нужно приложению, известно
    до первого шага.}
\end{frame}

% 9
\begin{frame}[fragile]{\texttt{Build} проверяет}
  \slidetop
  \claim{В Development контейнер проверяет \hi{каждую} регистрацию.}

\begin{lstlisting}[style=csharp,basicstyle=\codesm]
builder.Services.AddScoped<OrderService>();   // IClock is not registered
var app = builder.Build();                    // Development: throws here

sealed class OrderService(IClock clock) { }
\end{lstlisting}

\begin{lstlisting}[style=term,basicstyle=\termsm]
Unhandled exception. System.AggregateException: Some services are not able to be constructed ...
\end{lstlisting}

  \slidegap
  {\small Дальше в той же строке: \texttt{\hi{Unable to resolve service for type
    'IClock' while attempting to activate 'OrderService'.}}\par}

  \slidegap
  \aside{Падать \hi{сразу и громко} --- на \texttt{Build}, а не на первом
    запросе через неделю после выкладки. В Production проверки нет: \texttt{Build}
    проходит.}
\end{frame}

% 10
\begin{frame}[fragile]{Маршрут: путь и лямбда}
  \slidetop
  \claim{Обработчик --- это \hi{делегат}. Ответ из него делает фреймворк.}

  \begin{lstlisting}[style=csharp,basicstyle=\codexs]
app.MapGet("/", () => "Hello from BSUIR");
app.MapGet("/order", () => new { id = 42, status = "paid" });
\end{lstlisting}

  \begin{minipage}[t]{0.48\textwidth}
\begin{lstlisting}[style=term,basicstyle=\termsm]
$ curl -i http://localhost:5021/
HTTP/1.1 200 OK
Content-Type: text/plain; charset=utf-8
Date: Tue, 03 Nov 2026 12:08:29 GMT
Server: Kestrel
Transfer-Encoding: chunked

Hello from BSUIR
\end{lstlisting}
  \end{minipage}\hfill
  \begin{minipage}[t]{0.48\textwidth}
\begin{lstlisting}[style=term,basicstyle=\termsm]
$ curl -i http://localhost:5021/order
HTTP/1.1 200 OK
Content-Type: application/json; charset=utf-8
Date: Tue, 03 Nov 2026 12:13:46 GMT
Server: Kestrel
Transfer-Encoding: chunked

{"id":42,"status":"paid"}
\end{lstlisting}
  \end{minipage}

  \slidegap
  \aside{Строка уходит текстом, объект --- JSON. Минимальный API подробно ---
    в следующей лекции.}
\end{frame}

% 11
\begin{frame}{\texttt{Run}}
  \slidetop
  \claim{До \texttt{Run} --- один раз. Конвейер --- на \hi{каждый} запрос.}

  \begin{center}
    \begin{tikzpicture}[font=\footnotesize]
      \node[box,minimum width=5.2cm,minimum height=2.3cm,align=left] (st) at (0,0)
        {\texttt{CreateBuilder}\\\texttt{Services.Add...}\\\texttt{Build}\\
         \texttt{Use...} / \texttt{Map...}};
      \node[font=\scriptsize,text=muted,above=0.2em] at (st.north) {старт: один раз};
      \node[box,draw=accent,line width=1.0pt,minimum width=1.7cm] (run) at (4.4,0)
        {\texttt{Run}};
      \node[box,minimum width=4.4cm,minimum height=2.3cm] (rq) at (8.9,0)
        {конвейер\\на каждый запрос};
      \node[font=\scriptsize,text=accent,above=0.2em] at (rq.north)
        {пока не остановят};
      \draw[arr] (st) -- (run);
      \draw[arra] (run) -- (rq);
    \end{tikzpicture}
  \end{center}

  \slidegap
  \aside{\texttt{Run} не возвращается, пока приложение не остановят.}
\end{frame}

% 12
\begin{frame}{Где ломаться}
  \slidetop
  \versus{на старте}{%
    \begin{cbul}{accent}
      \item процесс не встал --- видно сразу
      \item старая версия продолжает работать
      \item клиенты ничего не заметили
    \end{cbul}}{на первом запросе}{%
    \begin{cbul}{alt}
      \item новая версия уже под нагрузкой
      \item падает она --- при клиентах
      \item через неделю, в три часа ночи
    \end{cbul}}

  \slidegap
  \hrulethin{0.42}
  \rulegap
  {\large Всё, что может быть настроено неправильно, ломается \hi{до Run}.}
\end{frame}

% 13
\begin{frame}[fragile]{Две фазы}
  \slidetop
  \begin{center}
    \begin{tikzpicture}[font=\footnotesize]
      \node[box,minimum width=5.4cm,minimum height=1.9cm,align=left] (l) at (0,0)
        {конфигурация\\сервисы\\логирование\\сервер};
      \node[box,draw=accent,minimum width=5.4cm,minimum height=1.9cm,align=left]
        (r) at (7.6,0) {middleware\\маршруты};
      \node[font=\scriptsize\bfseries,text=muted,above=0.3em] at (l.north)
        {ДО BUILD --- ОПИСАНИЕ};
      \node[font=\scriptsize\bfseries,text=accent,above=0.3em] at (r.north)
        {ПОСЛЕ BUILD --- КОНВЕЙЕР};
      \draw[alt,line width=1.4pt] (3.8,-1.15) -- (3.8,1.15);
      \node[font=\small\bfseries\ttfamily,fill=paper,text=alt] at (3.8,0) {Build};
    \end{tikzpicture}
  \end{center}

  \slidegap
{\small Сервис после \texttt{Build}: \texttt{\lo{The service collection cannot be
    modified because it is read-only.}}\par}

  \slidegap
  \aside{Две фазы \hi{не смешиваются}: глядя на файл, видно сначала, из чего
    приложение состоит, потом --- как оно обрабатывает запрос.}
\end{frame}

% 14
\begin{frame}{Хост}
  \slidetop
  \claim{Хост --- контейнер для всего, что живёт в \hi{процессе}.}

  \begin{center}
    \begin{tikzpicture}[font=\footnotesize,
      hb/.style={box,minimum width=2.3cm,minimum height=1.0cm}]
      \node[hb] (c1) at (0,0)    {конфигурация};
      \node[hb] (c2) at (2.6,0)  {контейнер};
      \node[hb] (c3) at (5.2,0)  {логгер};
      \node[hb,draw=accent,line width=1.0pt] (c4) at (7.8,0) {сервер\\Kestrel};
      \node[hb] (c5) at (10.4,0) {фоновые\\сервисы};
      \node[draw=muted,line width=0.9pt,rounded corners=3pt,inner sep=9pt,
            fit=(c1)(c5)] (h) {};
      \node[anchor=south west,font=\scriptsize\bfseries,text=muted]
        at ([yshift=2pt]h.north west) {ХОСТ: ЗАПУСК, ОСТАНОВКА, СОСТОЯНИЕ};
    \end{tikzpicture}
  \end{center}

  \slidegap
  \begin{cbul}{accent}
    \item сервер --- только \hi{один} из жильцов, самый заметный
    \item без сервера хост тоже есть: консольная утилита, обработчик очереди
    \item конфигурация, контейнер, логгер --- одинаковы в любом приложении
  \end{cbul}
\end{frame}

% 15
\begin{frame}{Три builder}
  \slidetop
  \claim{Один хост, три \hi{набора} умолчаний.}

  {\renewcommand{\arraystretch}{1.6}
  \begin{tabular}{@{}L{6.6cm}L{6.2cm}@{}}
    \hi{\texttt{WebApplication.CreateBuilder}}     & веб-приложение: всё, что
      было сегодня \\
    \texttt{Host.CreateApplicationBuilder}          & рабочий сервис: те же фазы,
      без порта (\texttt{dotnet new worker}) \\
    \texttt{WebApplication.CreateSlimBuilder}       & меньше умолчаний --- для
      обрезки и AOT \\
  \end{tabular}}

  \slidegap
  \aside{Те же \texttt{Services}, тот же \texttt{Build} и \texttt{Run}. Slim ---
    тот же хост, только диетический.}
\end{frame}

% 16
\begin{frame}{Что происходит на \texttt{Run}}
  \slidetop
  \begin{center}
    \begin{tikzpicture}[font=\footnotesize,
      sb/.style={box,minimum width=2.6cm,minimum height=1.2cm}]
      \node[sb] (s1) at (0,0)    {фоновые\\сервисы стартуют};
      \node[sb] (s2) at (3.3,0)  {сервер\\слушает порт};
      \node[sb,draw=accent,line width=1.0pt] (s3) at (6.6,0)
        {«приложение\\запущено»};
      \node[sb] (s4) at (9.9,0)  {ожидание\\сигнала};
      \draw[arr] (s1) -- (s2);
      \draw[arr] (s2) -- (s3);
      \draw[arr] (s3) -- (s4);
      \node[font=\scriptsize,text=muted,below=0.3em] at (s1.south) {по порядку регистрации};
      \node[font=\scriptsize,text=muted,below=0.3em] at (s3.south) {подписчики узнают};
    \end{tikzpicture}
  \end{center}

  \slidegap
  \aside{Дальше хост ничего не делает сам: сервер обрабатывает запросы,
    фоновые сервисы --- свою работу, хост ждёт остановки.}
\end{frame}

% 17
\begin{frame}[plain]
  \vote{Откуда приходит сигнал остановки?}
\end{frame}

% 18
\begin{frame}{Корректная остановка}
  \slidetop
  \begin{minipage}[t]{0.47\textwidth}
    \begin{tikzpicture}[font=\footnotesize,baseline=(current bounding box.north),
      st/.style={box,minimum width=5.6cm,minimum height=0.62cm,inner sep=2pt}]
      \node[st,draw=alt,line width=1.0pt] (t1) at (0,0)     {сигнал «завершайся»};
      \node[st] (t2) at (0,-0.9)  {сервер не принимает новые соединения};
      \node[st,draw=accent] (t3) at (0,-1.8) {текущие запросы дорабатывают};
      \node[st,draw=accent] (t4) at (0,-2.7) {фоновые сервисы получают токен};
      \node[st] (t5) at (0,-3.6)  {контейнер вызывает \texttt{Dispose}};
      \node[st] (t6) at (0,-4.5)  {процесс завершён};
      \foreach \a/\b in {t1/t2,t2/t3,t3/t4,t4/t5,t5/t6} \draw[arr] (\a) -- (\b);
    \end{tikzpicture}
  \end{minipage}\hfill
  \begin{minipage}[t]{0.47\textwidth}
    \vspace{0pt}
    \begin{cbul}{accent}
      \item локально --- Ctrl+C
      \item на сервере --- сигнал, и время, обычно \hi{десятки секунд}
      \item потом --- \lo{убивают}
      \item клиент, приславший запрос за секунду до сигнала, получит
        \hi{ответ}, а не ошибку
    \end{cbul}
  \end{minipage}
\end{frame}

% 19
\begin{frame}[fragile]{Таймаут остановки}
  \slidetop
  \claim{Таймаут хоста \hi{меньше}, чем ждёт тот, кто запустил процесс.}

\begin{lstlisting}[style=csharp,basicstyle=\codesm]
builder.Services.Configure<HostOptions>(options =>
    options.ShutdownTimeout = TimeSpan.FromSeconds(20));
\end{lstlisting}

  \slidegap
  {\renewcommand{\arraystretch}{1.45}
  \begin{tabular}{@{}L{4.2cm}L{8.6cm}@{}}
    по умолчанию      & \hi{30 секунд} \\
    не уложились      & работа \lo{брошена}: процесс завершается, незаконченное
      теряется \\
    фоновая работа на минуту & прерывается по токену за секунды, доделывает в
      следующий раз \\
  \end{tabular}}
\end{frame}

% 20
\begin{frame}{\texttt{Dispose} при остановке}
  \slidetop
  \claim{Контейнер освобождает то, чем \hi{владеет}, в обратном порядке.}

  \begin{center}
    \begin{tikzpicture}[font=\footnotesize,
      db/.style={box,minimum width=2.8cm,minimum height=0.8cm}]
      \node[db] (a) at (0,0)   {соединение с БД};
      \node[db] (b) at (4.0,0) {\texttt{LogWriter}};
      \node[db] (c) at (8.0,0) {кэш};
      \draw[arr] ([yshift=0.2cm]a.east) -- ([yshift=0.2cm]b.west);
      \draw[arr] ([yshift=0.2cm]b.east) -- ([yshift=0.2cm]c.west);
      \draw[arra] ([yshift=-0.2cm]c.west) -- ([yshift=-0.2cm]b.east);
      \draw[arra] ([yshift=-0.2cm]b.west) -- ([yshift=-0.2cm]a.east);
      \node[font=\scriptsize,text=muted,anchor=west] at (9.6,0.2) {создание};
      \node[font=\scriptsize,text=accent,anchor=west] at (9.6,-0.2) {\texttt{Dispose}};
    \end{tikzpicture}
  \end{center}

  \slidegap
  \begin{cbul}{accent}
    \item соединения закрыты, файлы и буферы \hi{сброшены}
    \item \texttt{LogWriter} дописывает буфер --- здесь
    \item созданное через \texttt{new} --- \lo{ваше}: контейнер о нём не знает
  \end{cbul}
\end{frame}

% 21
\begin{frame}{Не обходите остановку}
  \slidetop
  \versus{корректно}{%
    \begin{cbul}{accent}
      \item сигнал или \texttt{StopApplication}
      \item цикл смотрит на токен
      \item \texttt{Dispose} всего, чем владеет хост
    \end{cbul}}{в обход}{%
    \begin{cbul}{alt}
      \item \texttt{Environment.Exit} из кода
      \item токен игнорируется
      \item процесс убит по таймауту
    \end{cbul}}

  \slidegap
  \hrulethin{0.42}
  \rulegap
  {\large Цена: обрезанные ответы, недописанные файлы, \lo{незакрытые}
    транзакции.}
\end{frame}

% 22
\begin{frame}[fragile]{Жизненный цикл из кода}
  \slidetop
  \claim{На события хоста можно \hi{подписаться}.}

\begin{lstlisting}[style=csharp,basicstyle=\codesm]
var lifetime = app.Services.GetRequiredService<IHostApplicationLifetime>();
lifetime.ApplicationStarted.Register(() => app.Logger.LogInformation("Started"));
lifetime.ApplicationStopping.Register(() => app.Logger.LogInformation("Stopping"));

app.MapPost("/shutdown", (IHostApplicationLifetime l) =>
{
    l.StopApplication();
    return Results.Accepted();
});
\end{lstlisting}

  \slidegap
  \aside{\texttt{StopApplication} --- остановка тем же корректным путём, а
    не убийство процесса.}
\end{frame}

% 23
\begin{frame}[fragile]{Первые строки лога}
  \slidetop
  \claim{Половина «почему не открывается» --- \hi{здесь}.}

\begin{lstlisting}[style=term,basicstyle=\termsm]
$ dotnet run
Using launch settings from /home/user/HelloWorld/Properties/launchSettings.json...
Building...
info: Microsoft.Hosting.Lifetime[14]
      Now listening on: http://localhost:5012
info: Microsoft.Hosting.Lifetime[0]
      Application started. Press Ctrl+C to shut down.
info: Microsoft.Hosting.Lifetime[0]
      Hosting environment: Development
info: Microsoft.Hosting.Lifetime[0]
      Content root path: /home/user/HelloWorld
^C
info: Microsoft.Hosting.Lifetime[0]
      Application is shutting down...
\end{lstlisting}

  \slidegap
  \aside{Слушает не там, окружение не то, папка не та --- видно в первых
    четырёх строках.}
\end{frame}

% 24
\begin{frame}{Окружение}
  \slidetop
  \versus{Development}{%
    \begin{cbul}{accent}
      \item страница с трассировкой при ошибке
      \item подробные логи
      \item проверка контейнера при \texttt{Build}
    \end{cbul}}{Production}{%
    \begin{cbul}{alt}
      \item наружу --- только статус
      \item логи по настройкам
      \item без проверки при \texttt{Build}
    \end{cbul}}

  \slidegap
  \hrulethin{0.42}
  \rulegap
  {\large Разница --- в \hi{конфигурации}, а не в коде.}

  \vspace{0.6em}
  \aside{Каждое «если разработка, то\ldots» в файле запуска --- развилка, которую
    придётся поддерживать.}
\end{frame}

% 25
\begin{frame}[fragile]{Чего не видно в \texttt{Program.cs}}
  \slidetop
  \begin{minipage}[t]{0.46\textwidth}
    \eyebrow{HelloWorld.csproj}\par\vspace{0.3em}
\begin{lstlisting}[style=xml,basicstyle=\termsm]
<Project Sdk="Microsoft.NET.Sdk.Web">

  <PropertyGroup>
    <TargetFramework>net10.0</TargetFramework>
    <Nullable>enable</Nullable>
    <ImplicitUsings>enable</ImplicitUsings>
  </PropertyGroup>

</Project>
\end{lstlisting}
    \vspace{0.4em}
    {\small\raggedright \hi{Sdk.Web}: веб-часть рантайма, неявные \texttt{using},
      файлы настроек в выводе сборки.\par}
  \end{minipage}\hfill
  \begin{minipage}[t]{0.50\textwidth}
    \eyebrow{Properties / launchSettings.json}\par\vspace{0.3em}
\begin{lstlisting}[style=json,basicstyle=\termsm]
"http": {
  "commandName": "Project",
  "dotnetRunMessages": true,
  "launchBrowser": true,
  "applicationUrl": "http://localhost:5021",
  "environmentVariables": {
    "ASPNETCORE_ENVIRONMENT": "Development"
  }
}
\end{lstlisting}
    \vspace{0.4em}
    {\small\raggedright Читает \texttt{dotnet run}: адрес и окружение.
      \lo{Только локально} --- в прод не едет.\par}
  \end{minipage}
\end{frame}

% 26
\begin{frame}[plain]
  \keyline{Приложение --- это хост плюс конвейер.\\ До \texttt{Build}
    описываем, после --- собираем конвейер, \texttt{Run} запускает и корректно
    останавливает.}
\end{frame}

% ======================================================== 2. KESTREL ==
\section{2. Kestrel и адрес}

% 27
\begin{frame}[plain]
  \sectionopener{2}{Kestrel и адрес}{Kestrel --- это ваш сервер на сокетах,
    только \hi{правильный}.}
\end{frame}

% 28
\begin{frame}{Kestrel и сервер из лабораторной}
  \slidetop
  \versus{Kestrel}{%
    \begin{cbul}{accent}
      \item HTTP/1.1, HTTP/2 и HTTP/3
      \item TLS
      \item пулы буферов
      \item лимиты на размер и время
    \end{cbul}}{сервер на сокетах}{%
    \begin{cbul}{alt}
      \item только HTTP/1.1
      \item без TLS
      \item буфер на каждый запрос
      \item без лимитов
    \end{cbul}}

  \slidegap
  \hrulethin{0.42}
  \rulegap
  {\large Делает то же самое: принимает, разбирает, отвечает. Вы к нему
    \hi{не прикасаетесь}.}
\end{frame}

% 29
\begin{frame}[fragile]{Где слушать}
  \slidetop
  \claim{Адрес --- свойство \hi{окружения}, а не приложения.}

  \begin{lstlisting}[style=term,basicstyle=\termsm]
$ dotnet HelloWorld.dll
      Now listening on: http://localhost:5000
$ ASPNETCORE_URLS=http://localhost:8092 dotnet HelloWorld.dll
      Now listening on: http://localhost:8092
$ ASPNETCORE_URLS=http://localhost:8092 dotnet HelloWorld.dll --urls=http://localhost:8091
      Now listening on: http://localhost:8091
$ ASPNETCORE_HTTP_PORTS=8093 dotnet HelloWorld.dll
      Now listening on: http://[::]:8093
\end{lstlisting}

  \slidegap
  {\renewcommand{\arraystretch}{1.3}\small
  \begin{tabular}{@{}L{3.2cm}L{9.6cm}@{}}
    \hi{в контейнере}   & все интерфейсы, порт 8080 --- его задаёт образ \\
    \lo{секция Kestrel} & в \texttt{appsettings.json} перекрывает \texttt{-{}-urls} и
      \texttt{ASPNETCORE\_URLS} --- выбирайте \lo{один} способ \\
  \end{tabular}}
\end{frame}

% 30
\begin{frame}[fragile]{HTTPS локально}
  \slidetop
  \claim{Сертификат для разработки создаёт \hi{SDK}.}

\begin{lstlisting}[style=bash,basicstyle=\ttfamily\small]
dotnet dev-certs https --trust
\end{lstlisting}

  \slidegap
  \begin{cbul}{accent}
    \item сертификат для \texttt{localhost}, подписан \lo{сам собой}
    \item браузер не верит, пока не добавите в доверенные --- это и делает
      \texttt{-{}-trust}
    \item в проде сертификат --- \hi{конфигурация}, а не код
  \end{cbul}
\end{frame}

% 31
\begin{frame}{Версии протокола}
  \slidetop
  \claim{Ваш код от версии \hi{не зависит}.}

  {\renewcommand{\arraystretch}{1.6}
  \begin{tabular}{@{}L{2.6cm}L{10.2cm}@{}}
    \hi{\textbf{HTTP/1.1}} & по умолчанию \\
    \hi{\textbf{HTTP/2}}   & по умолчанию; браузер договаривается о нём только
      поверх TLS \\
    \textbf{HTTP/3}        & включается явно; нужна поддержка QUIC в ОС \\
  \end{tabular}}

  \slidegap
  \aside{Семантика одна для всех трёх версий --- транспорт под ней меняется,
    обработчик нет.}
\end{frame}

% 32
\begin{frame}{Лимиты по умолчанию}
  \slidetop
  \claim{Лишнее отсекает \hi{сервер}, а не ваш код.}

  {\renewcommand{\arraystretch}{1.25}\small
  \begin{tabular}{@{}L{5.2cm}L{2.9cm}L{4.7cm}@{}}
    \colhead{muted}{лимит} & \colhead{muted}{по умолчанию} & \colhead{muted}{если больше} \\
    \texttt{MaxRequestBodySize}         & \hi{30\,000\,000 байт} & \lo{413} при чтении тела \\
    \texttt{MaxRequestHeadersTotalSize} & 32 КБ  & отказ до вашего кода \\
    \texttt{MaxRequestHeaderCount}      & 100    & отказ до вашего кода \\
    \texttt{RequestHeadersTimeout}      & 30 с   & медленный клиент отключён \\
    \texttt{KeepAliveTimeout}           & 2 мин 10 с & соединение закрыто \\
    \texttt{MaxConcurrentConnections}   & без предела & --- \\
  \end{tabular}}
\end{frame}

% 33
\begin{frame}[fragile]{Поднять лимит}
  \slidetop
  \claim{Для всего сервера или для \hi{одного} маршрута.}

  \begin{lstlisting}[style=csharp,basicstyle=\codexs]
builder.WebHost.ConfigureKestrel(kestrel =>
    kestrel.Limits.MaxRequestBodySize = 100_000_000);

app.MapPost("/upload", [RequestSizeLimit(100_000_000)] async (HttpRequest request) =>
{
    await request.Body.CopyToAsync(Stream.Null);
    return Results.Ok();
});
\end{lstlisting}

  \slidegap
  {\small Лимит по умолчанию, 31 МБ в тело --- отказ на \hi{первом чтении}:\par}
\begin{lstlisting}[style=term,basicstyle=\termsm]
$ curl -i --data-binary @big.bin http://localhost:5021/upload
HTTP/1.1 413 Payload Too Large
Connection: close
\end{lstlisting}
\end{frame}

% 34
\begin{frame}[plain]
  \keyline{Kestrel отсекает лишнее за вас.\\ Где слушать и по какой схеме ---
    решает окружение.}
\end{frame}

% ======================================================== 3. PIPELINE ==
\section{3. Конвейер запроса}

% The chain of six boxes, drawn the same way on frames 37 and 59.
\newcommand{\chainboxes}{%
  \foreach \i/\t in {1/ошибки,2/логирование,3/маршрутизация,4/аутентификация,
                     5/авторизация} {
    \node[box,minimum width=1.95cm,minimum height=1.0cm,inner sep=1pt,font=\scriptsize]
      (m\i) at ({(\i-1)*2.32},0) {\t};
  }
  \node[box,draw=accent,line width=1.0pt,minimum width=1.95cm,minimum height=1.0cm,
        inner sep=1pt,font=\scriptsize\bfseries] (m6) at (11.6,0) {обработчик};
}

% 35
\begin{frame}[plain]
  \sectionopener{3}{Конвейер запроса}{Middleware --- это делегат, который
    \hi{знает следующего}.}
\end{frame}

% 36
\begin{frame}{Простая модель ломается}
  \slidetop
  \claim{Каждому запросу нужно \hi{одно и то же}.}

  \begin{minipage}[t]{0.40\textwidth}
    \begin{cbul}{accent}
      \item записать в лог
      \item узнать, кто прислал
      \item решить, можно ли
      \item сжать ответ
      \item поймать исключение
      \item отдать файл
    \end{cbul}
  \end{minipage}\hfill
  \begin{minipage}[t]{0.54\textwidth}
    \vspace{0pt}
    {\renewcommand{\arraystretch}{1.6}
    \begin{tabular}{@{}L{7.3cm}@{}}
      \lo{в каждом обработчике} --- сто копий одних и тех же двадцати строк \\
      \lo{в одном месте перед обработчиком} --- метод на тысячу строк с
        флагами \\
    \end{tabular}}
  \end{minipage}

  \slidegap
  \aside{Нужен \hi{третий} вариант.}
\end{frame}

% 37
\begin{frame}{Конвейер}
  \slidetop
  \claim{Каждый компонент видит запрос \hi{до} обработчика и ответ \hi{после}.}

  \begin{center}
    \begin{tikzpicture}
      \chainboxes
      \foreach \i [evaluate=\i as \j using int(\i+1)] in {1,...,5} {
        \draw[arra] ([yshift=0.22cm]m\i.east) -- ([yshift=0.22cm]m\j.west);
        \draw[arrl] ([yshift=-0.22cm]m\j.west) -- ([yshift=-0.22cm]m\i.east);
      }
      \node[anchor=south west,font=\scriptsize,text=accent,inner sep=1pt]
        at ([yshift=2pt]m1.north west) {запрос $\to$};
      \node[anchor=north west,font=\scriptsize,text=alt,inner sep=1pt]
        at ([yshift=-2pt]m1.south west) {$\leftarrow$ ответ};
    \end{tikzpicture}
  \end{center}

  \slidegap
  \begin{cbul}{accent}
    \item каждый делает \hi{одно} дело и передаёт дальше
    \item добавить компонент --- вставить ещё одно звено, остальные не меняются
    \item ответ идёт по той же цепочке \hi{обратно}
  \end{cbul}
\end{frame}

% 38
\begin{frame}[plain]
  \statement{\hi{Middleware} --- делегат, который получает контекст и
    следующего и сам решает: вызвать ли следующего, что сделать до и что
    после.}
\end{frame}

% 39
\begin{frame}[fragile]{Явный конвейер в \texttt{Program.cs}}
  \vspace{0.3em}
\begin{lstlisting}[style=csharp,basicstyle=\ttfamily\fontsize{7.0}{7.9}\selectfont,escapechar=|]
var builder = WebApplication.CreateBuilder(args);

builder.Services.AddAuthentication();
builder.Services.AddAuthorization();
builder.Services.AddHealthChecks();

var app = builder.Build();

|\mk{e1}|app.UseMiddleware<ErrorHandlingMiddleware>();
|\mk{e2}|app.UseMiddleware<RequestLoggingMiddleware>();
|\mk{e3}|app.UseRouting();
|\mk{e4}|app.UseAuthentication();
|\mk{e5}|app.UseAuthorization();

|\mk{e6}|app.MapHealthChecks("/health");
|\mk{e7}|app.MapGet("/", () => "Hello from BSUIR");

app.Run();
\end{lstlisting}
  \begin{tikzpicture}[remember picture,overlay,every node/.style={rawlab}]
    \node at ([xshift=7.6cm]e1) {\hi{\textbf{1}} вокруг всего: ошибки};
    \node at ([xshift=7.6cm]e2) {\hi{\textbf{2}} каждый запрос в лог};
    \node at ([xshift=7.6cm]e3) {\hi{\textbf{3}} выбрать обработчик};
    \node at ([xshift=7.6cm]e4) {\hi{\textbf{4}} кто пришёл};
    \node at ([xshift=7.6cm]e5) {\hi{\textbf{5}} можно ли};
    \draw[rawbrace] ([xshift=7.35cm,yshift=0.35em]e6) -- ([xshift=7.35cm,yshift=-0.3em]e7);
    \node[text=accent] at ([xshift=7.6cm]$(e6)!0.5!(e7)$) {конец цепочки: обработчики};
  \end{tikzpicture}

  \vspace{-0.1em}
  \aside{Сверху вниз идёт \hi{запрос}, снизу вверх --- ответ.}
\end{frame}

% 40
\begin{frame}{Три возможности}
  \slidetop
  \claim{До, вызвать или \hi{не} вызвать, после --- из этого собирается всё.}

  {\renewcommand{\arraystretch}{1.6}
  \begin{tabular}{@{}L{4.2cm}L{8.6cm}@{}}
    \hi{\textbf{до и после}}   & логирование: засечь время, записать статус \\
    \hi{\textbf{до}}           & аутентификация: узнать, кто пришёл \\
    \hi{\textbf{вокруг}}       & обработка ошибок: \texttt{try} вокруг всего
      остального \\
    \lo{\textbf{не вызывать}}  & статический файл: ответить самому \\
  \end{tabular}}
\end{frame}

% 41
\begin{frame}[fragile]{Матрёшка}
  \slidetop
  \claim{При старте цепочка складывается в \hi{один} делегат.}

  \begin{minipage}[t]{0.50\textwidth}
\begin{lstlisting}[style=csharp,basicstyle=\codesm]
// one request
delegate Task RequestDelegate(
    HttpContext context);

// one middleware
Func<RequestDelegate, RequestDelegate>
\end{lstlisting}
  \end{minipage}\hfill
  \begin{minipage}[t]{0.45\textwidth}
    \vspace{0.3em}
    \begin{tikzpicture}[font=\scriptsize]
      \node[draw=hairline,line width=0.8pt,minimum width=5.6cm,minimum height=3.0cm,
            label={[anchor=north west,text=muted]north west:ошибки}] at (0,0) {};
      \node[draw=hairline,line width=0.8pt,minimum width=4.6cm,minimum height=2.2cm,
            label={[anchor=north west,text=muted]north west:логирование}] at (0,-0.2) {};
      \node[draw=hairline,line width=0.8pt,minimum width=3.6cm,minimum height=1.4cm,
            label={[anchor=north west,text=muted]north west:маршрутизация}] at (0,-0.4) {};
      \node[box,draw=accent,minimum width=2.4cm,minimum height=0.6cm,
            inner sep=1pt] at (0,-0.6) {обработчик};
    \end{tikzpicture}
  \end{minipage}

  \slidegap
  \aside{Последний оборачивается предпоследним, тот --- предыдущим. На запрос
    вызывается один делегат: \hi{без цикла}, без рефлексии, без поиска.}
\end{frame}

% 42
\begin{frame}[fragile]{Middleware одной лямбдой}
  \vspace{0.3em}
\begin{lstlisting}[style=csharp,basicstyle=\codesm,escapechar=|]
using System.Diagnostics;

app.Use(async (context, next) =>
{
|\mk{i1}|    var started = Stopwatch.GetTimestamp();

|\mk{i2}|    await next(context);

|\mk{i3}|    var elapsed = Stopwatch.GetElapsedTime(started);
    app.Logger.LogInformation("{Method} {Path} -> {Status} in {Elapsed} ms",
        context.Request.Method, context.Request.Path,
        context.Response.StatusCode, elapsed.TotalMilliseconds);
});
\end{lstlisting}
  \begin{tikzpicture}[remember picture,overlay,every node/.style={rawlab}]
    \node at ([xshift=9.0cm]i1) {\hi{до next} --- запрос};
    \node at ([xshift=9.0cm]i2) {\hi{await next} --- всё остальное};
    \node at ([xshift=9.0cm]i3) {\hi{после next} --- ответ};
  \end{tikzpicture}

  \vspace{0.2em}
  \aside{Шаблон с именами в фигурных скобках и параметры отдельно --- не
    интерполяция.}
\end{frame}

% 43
\begin{frame}[fragile]{После \texttt{next} --- поздно}
  \slidetop
  \claim{Первый байт тела ушёл --- заголовки \lo{не вернуть}.}

  \begin{lstlisting}[style=term,basicstyle=\termsm]
System.InvalidOperationException: StatusCode cannot be set because the response has already started.
\end{lstlisting}

  \slidegap
  {\renewcommand{\arraystretch}{1.4}
  \begin{tabular}{@{}L{3.6cm}L{9.2cm}@{}}
    \hi{до \texttt{next}}    & менять статус и заголовки --- можно \\
    \lo{после \texttt{next}} & читать статус, писать в лог --- можно; менять ---
      поздно \\
  \end{tabular}}
\end{frame}

% 44
\begin{frame}{Порядок --- это семантика}
  \slidetop
  \claim{Первый видит запрос \hi{первым}, а ответ --- \hi{последним}.}

  {\renewcommand{\arraystretch}{1.6}
  \begin{tabular}{@{}L{5.4cm}L{7.4cm}@{}}
    \hi{\textbf{аутентификация до авторизации}} & иначе решают про неизвестно
      кого \\
    \hi{\textbf{ошибки снаружи всех}}           & иначе исключение улетит в сервер:
      пустой 500 \\
    \hi{\textbf{маршрутизация до авторизации}}  & иначе не по чему решать:
      требования у маршрута \\
  \end{tabular}}

  \slidegap
  \aside{Анализатор этого не ловит. Защищённый маршрут падает на первом запросе:
    \texttt{contains authorization metadata, but a middleware was not found that
    supports authorization}.}
\end{frame}

% 45
\begin{frame}[fragile]{Ветвление}
  \slidetop
  \claim{Конвейер не обязан быть \hi{одной} прямой.}

\begin{lstlisting}[style=csharp,basicstyle=\codexs,escapechar=|]
app.Map("/admin", admin => admin.Run(context => context.Response.WriteAsync("admin")));

app.UseWhen(context => context.Request.Path.StartsWithSegments("/api"),
    api => api.Use(async (context, next) =>
    {
        context.Response.Headers["X-Api"] = "1";
        await next(context);
    }));
\end{lstlisting}

  \slidegap
  {\renewcommand{\arraystretch}{1.4}
  \begin{tabular}{@{}L{2.2cm}L{10.6cm}@{}}
    \hi{\texttt{Map}}     & своя цепочка по префиксу пути; в основную
      \hi{не возвращается} \\
    \hi{\texttt{UseWhen}} & вставка по условию; запрос \hi{возвращается} в
      основную \\
  \end{tabular}}
\end{frame}

% 46
\begin{frame}{Конвейер по умолчанию}
  \slidetop
  \claim{До вашего первого \texttt{Use} там \hi{уже} есть компоненты.}

  {\renewcommand{\arraystretch}{1.3}
  \begin{tabular}{@{}L{0.5cm}L{5.4cm}L{6.9cm}@{}}
    \textcolor{muted}{1} & страница ошибок разработчика & \lo{только в Development}:
      трассировка наружу --- утечка \\
    \textcolor{muted}{2} & фильтрация хостов & чужой заголовок \texttt{Host} --- отказ \\
    \textcolor{muted}{3} & маршрутизация & выбирает обработчик, но \hi{не вызывает} \\
    \textcolor{muted}{4} & аутентификация и авторизация & если зарегистрированы ---
      после маршрутизации \\
    \textcolor{accent}{5} & \hi{\textbf{ваши middleware}} & там, где вы их написали \\
    \textcolor{muted}{6} & конечные точки & вызывают выбранный обработчик \\
  \end{tabular}}

  \rulegap
  \aside{Стандартные компоненты фреймворк ставит \hi{сам} и в правильном
    порядке.}
\end{frame}

% 47
\begin{frame}[fragile]{Как убедиться}
  \slidetop
  \claim{Один и тот же middleware \hi{до} маршрутизации и \hi{после}.}

  \begin{lstlisting}[style=csharp,basicstyle=\codexs]
app.Use(async (context, next) =>
{
    app.Logger.LogInformation("before UseRouting: {Endpoint}",
        context.GetEndpoint()?.DisplayName ?? "(none)");
    await next(context);
});
app.UseRouting();
// the same middleware again: "after UseRouting: {Endpoint}"
\end{lstlisting}

\begin{lstlisting}[style=term,basicstyle=\termsm]
info: HelloWorld[0]
      before UseRouting: (none)
info: HelloWorld[0]
      after UseRouting: HTTP: GET /
\end{lstlisting}

  \rulegap
  \aside{Без явного \texttt{UseRouting} маршрутизация стоит \hi{перед} вашим
    первым \texttt{Use}: он уже видит обработчик.}
\end{frame}

% 48
\begin{frame}{Свой middleware: три способа}
  \slidetop
  {\renewcommand{\arraystretch}{1.55}\small
  \begin{tabular}{@{}L{3.3cm}L{3.0cm}L{3.0cm}L{3.5cm}@{}}
    & \colhead{muted}{создаётся} & \colhead{muted}{зависимости} & \colhead{muted}{когда брать} \\
    \textbf{лямбда в \texttt{Use}}  & один раз & то, что видно в файле & две-три
      строки \\
    \hi{\textbf{класс с \texttt{InvokeAsync}}} & \hi{один раз} при старте &
      параметры \texttt{InvokeAsync} & \hi{всё остальное} \\
    \textbf{\texttt{IMiddleware}}   & на каждый запрос & конструктор & когда точно
      знаете, зачем \\
  \end{tabular}}

  \slidegap
  \aside{\texttt{IMiddleware} --- объект на каждый запрос. Это цена, и она нужна
    редко.}
\end{frame}

% 49
\begin{frame}[fragile]{Класс с \texttt{InvokeAsync}}
  \vspace{0.3em}
\begin{lstlisting}[style=csharp,basicstyle=\codesm,escapechar=|]
sealed class RequestTimingMiddleware(
|\mk{c1}|    RequestDelegate next,
|\mk{c2}|    ILogger<RequestTimingMiddleware> logger)
{
    public async Task InvokeAsync(
|\mk{c3}|        HttpContext context,
|\mk{c4}|        RequestStamp stamp)
    {
        var started = Stopwatch.GetTimestamp();
        await next(context);
        logger.LogInformation("{Path} [{Stamp}] in {Elapsed} ms",
            context.Request.Path, stamp.Value,
            Stopwatch.GetElapsedTime(started).TotalMilliseconds);
    }
}

// builder.Services.AddScoped<RequestStamp>();
\end{lstlisting}
  \begin{tikzpicture}[remember picture,overlay,every node/.style={rawlab}]
    \node at ([xshift=7.4cm]c1) {\hi{конструктор} --- один раз при старте};
    \node at ([xshift=7.4cm]c2) {только вечное: логгер, настройки};
    \node at ([xshift=7.4cm]c3) {\hi{InvokeAsync} --- на каждый запрос};
    \node at ([xshift=7.4cm]c4) {параметры --- из контейнера запроса};
  \end{tikzpicture}
\end{frame}

% 50
\begin{frame}[fragile]{Метод расширения \texttt{Use...}}
  \slidetop
  \claim{Файл запуска --- \hi{словарь}, а не список типов.}

\begin{lstlisting}[style=csharp,basicstyle=\codesm]
static class RequestTimingExtensions
{
    public static IApplicationBuilder UseRequestTiming(this IApplicationBuilder app)
        => app.UseMiddleware<RequestTimingMiddleware>();
}

app.UseRequestTiming();
\end{lstlisting}

  \slidegap
  \aside{Так устроены и стандартные: \texttt{UseRouting},
    \texttt{UseAuthentication}. Внутри --- одна строка.}
\end{frame}

% 51
\begin{frame}[fragile]{Пленённая зависимость}
  \slidetop
  \versus{параметры InvokeAsync}{%
    \begin{cbul}{accent}
      \item на каждый запрос --- свой экземпляр
      \item контекст базы, всё «на запрос»
    \end{cbul}}{конструктор}{%
    \begin{cbul}{alt}
      \item один раз --- \lo{навсегда}
      \item только то, что живёт вечно
    \end{cbul}}

  \slidegap
{\small В конструкторе, в Development --- падение на старте:\par
  \texttt{\lo{Cannot resolve scoped service 'RequestStamp' from root provider.}}\par}

  \rulegap
  \aside{В Production ошибки нет: два запроса получают \lo{один и тот же}
    экземпляр.}
\end{frame}

% 52
\begin{frame}[plain]
  \vote{Ловить всё можно только на границе процесса. Где эта граница в
    сервисе?}
\end{frame}

% 53
\begin{frame}[fragile]{Обработка ошибок --- middleware}
  \vspace{0.2em}
\begin{lstlisting}[style=csharp,basicstyle=\ttfamily\fontsize{6.6}{7.2}\selectfont,escapechar=|]
sealed class ErrorHandlingMiddleware(
    RequestDelegate next, ILogger<ErrorHandlingMiddleware> logger)
{
    public async Task InvokeAsync(HttpContext context)
    {
        try
        {
|\mk{h1}|            await next(context);
        }
|\mk{h2}|        catch (Exception e) when (!context.Response.HasStarted)
        {
|\mk{h3}|            var (status, message) = e switch
            {
                NotFoundException => (StatusCodes.Status404NotFound, e.Message),
                _ => (StatusCodes.Status500InternalServerError, "Internal error"),
            };
|\mk{h4}|            logger.LogError(e, "Request {Path} failed with {Status}",
                context.Request.Path, status);
            context.Response.StatusCode = status;
|\mk{h5}|            await context.Response.WriteAsJsonAsync(
                new { error = message, traceId = context.TraceIdentifier });
        }
    }
}
\end{lstlisting}
  \begin{tikzpicture}[remember picture,overlay,every node/.style={rawlab,font=\small\bfseries,text=accent}]
    \foreach \m/\n in {h1/1,h2/2,h3/3,h4/4,h5/5} \node at ([xshift=11.9cm]\m) {\n};
  \end{tikzpicture}
\end{frame}

% 54
\begin{frame}{Пять решений}
  \slidetop
  {\renewcommand{\arraystretch}{1.5}
  \begin{tabular}{@{}L{4.0cm}L{8.8cm}@{}}
    \hi{\textbf{1\enspace try вокруг next}}   & всё приложение внутри одного \texttt{try};
      поэтому он \hi{первый} \\
    \hi{\textbf{2\enspace when (!HasStarted)}} & ответ уже уходит --- не ловим, соединение
      обрывается \\
    \hi{\textbf{3\enspace switch по типу}}    & «не найдено» $\to$ 404, остальное $\to$ 500 \\
    \hi{\textbf{4\enspace в лог --- всё}}      & исключение целиком, первым параметром \\
    \lo{\textbf{5\enspace наружу --- минимум}} & для 500: «Internal error» и
      \texttt{traceId}, без сообщения исключения \\
  \end{tabular}}

  \slidegap
  \aside{Сообщение исключения может содержать имя таблицы, путь к файлу, кусок
    SQL. \texttt{NotFoundException} --- своё исключение домена, наследник
    \texttt{Exception}.}
\end{frame}

% 55
\begin{frame}[fragile]{Отмена --- не ошибка}
  \slidetop
  \claim{Клиент ушёл --- отвечать \hi{некому}.}

\begin{lstlisting}[style=csharp,basicstyle=\codexs]
try
{
    await next(context);
}
catch (OperationCanceledException) when (context.RequestAborted.IsCancellationRequested)
{
    logger.LogInformation("Request {Path} cancelled by the client", context.Request.Path);
}
catch (Exception e) when (!context.Response.HasStarted)
{
    // as before
}
\end{lstlisting}

  \slidegap
  \aside{Иначе лог забит «ошибками» о закрытых вкладках. Ветка --- в
    лабораторной.}
\end{frame}

% 56
\begin{frame}[fragile]{Проверка руками}
  \slidetop
  \begin{lstlisting}[style=term,basicstyle=\termsm]
$ curl -s http://localhost:5012/orders/404
{"error":"Order 404 not found","traceId":"0HNPPKMVENS8K:00000001"}
$ curl -i http://localhost:5012/orders/500
HTTP/1.1 500 Internal Server Error
Content-Type: application/json; charset=utf-8
Date: Tue, 03 Nov 2026 12:26:54 GMT
Server: Kestrel
Transfer-Encoding: chunked

{"error":"Internal error","traceId":"0HNPPKMVENS8L:00000001"}
\end{lstlisting}

  \slidegap
\begin{lstlisting}[style=term,basicstyle=\termsm]
fail: ErrorHandlingMiddleware[0]
      Request /orders/500 failed with 500
      System.InvalidOperationException: Connection failed: Host=db;Password=secret
\end{lstlisting}

  \rulegap
  \aside{Пароль --- \hi{только в логе}. Клиенту --- «Internal error» и
    \texttt{traceId}.}
\end{frame}

% 57
\begin{frame}{Одно место против ста}
  \slidetop
  \versus{один middleware}{%
    \begin{cbul}{accent}
      \item один формат ошибки
      \item решение принято и проверено один раз
      \item обработчики короткие: делают и бросают
    \end{cbul}}{try/catch в каждом}{%
    \begin{cbul}{alt}
      \item сто копий --- через год они разные
      \item где-то забыли --- голый 500
      \item где-то «для понятности» --- строка подключения наружу
    \end{cbul}}
\end{frame}

% 58
\begin{frame}{Принцип один}
  \slidetop
  \claim{Какой именно код это делает --- \hi{вторично}.}

  \begin{cbul}{accent}
    \item \hi{одно} место
    \item снаружи \hi{всех}
    \item наружу --- минимум
    \item в лог --- всё
  \end{cbul}

  \slidegap
  \aside{У фреймворка есть встроенный механизм с \texttt{IExceptionHandler} и
    стандартным форматом тела. Мы пишем свой, чтобы видеть каждую строку.}
\end{frame}

% 59
\begin{frame}{Короткое замыкание}
  \slidetop
  \claim{Не вызвал \texttt{next} --- запрос \hi{дальше не идёт}.}

  \begin{center}
    \begin{tikzpicture}
      \chainboxes
      \draw[arra] ([yshift=0.22cm]m1.east) -- ([yshift=0.22cm]m2.west);
      \draw[arra] ([yshift=0.22cm]m2.east) -- ([yshift=0.22cm]m3.west);
      \draw[arrl] ([yshift=-0.22cm]m3.west) -- ([yshift=-0.22cm]m2.east);
      \draw[arrl] ([yshift=-0.22cm]m2.west) -- ([yshift=-0.22cm]m1.east);
      \node[draw=accent,line width=1.2pt,fit=(m3),inner sep=1pt] {};
      \foreach \i in {4,5,6} \node[fill=paper,opacity=0.6,fit=(m\i),inner sep=1pt] {};
      \node[anchor=south west,font=\scriptsize,text=accent,inner sep=1pt]
        at ([yshift=2pt]m1.north west) {запрос $\to$};
      \node[font=\footnotesize,text=accent,below=0.55cm] at (m3.south)
        {ответ сформирован здесь};
    \end{tikzpicture}
  \end{center}

  \slidegap
  \aside{Это не ошибка, а \hi{инструмент}: middleware сам пишет ответ, и цепочка
    разворачивается на нём.}
\end{frame}

% 60
\begin{frame}{Кто отвечает сам}
  \slidetop
  {\renewcommand{\arraystretch}{1.55}
  \begin{tabular}{@{}L{4.6cm}L{2.0cm}L{6.2cm}@{}}
    \colhead{muted}{middleware} & \colhead{muted}{ответ} & \colhead{muted}{почему не дальше} \\
    проверка здоровья          & \hi{200}  & вашему коду там нечего делать \\
    статические файлы          & файл      & обработчики не знают, что файлы есть \\
    ограничение частоты        & \lo{429}  & нагрузка на базу не растёт \\
    предварительный \texttt{OPTIONS} & 204 & обработчику там нечего делать \\
    аутентификация, нет токена & \lo{401}  & обработчик не вызван \\
  \end{tabular}}

  \slidegap
  \aside{Срабатывает до всего, что стоит после. Стоит \lo{после} проверки
    токена --- сначала потребует токен.}
\end{frame}

% 61
\begin{frame}{Middleware и фильтр}
  \slidetop
  \versus{middleware}{%
    \begin{cbul}{accent}
      \item все запросы
      \item про обработчик не знает
      \item лог, ошибки, сжатие
    \end{cbul}}{фильтр конечной точки}{%
    \begin{cbul}{alt}
      \item один маршрут
      \item видит параметры обработчика
      \item валидация входа
    \end{cbul}}

  \slidegap
  \hrulethin{0.42}
  \rulegap
  {\large Одинаково для всех --- \hi{middleware}. Зависит от маршрута ---
    фильтр.}
\end{frame}

% 62
\begin{frame}[fragile]{Забытый \texttt{next}}
  \slidetop
  \claim{Все запросы обрываются на \lo{нём}.}

  \begin{lstlisting}[style=csharp,basicstyle=\codexs]
app.Use(async (HttpContext context, RequestDelegate next) =>
{
    app.Logger.LogInformation("Request {Path}", context.Request.Path);
    // forgot: await next(context);
});
\end{lstlisting}

\begin{lstlisting}[style=term,basicstyle=\termsm]
$ curl -i http://localhost:5021/anything
HTTP/1.1 200 OK
Content-Length: 0
Date: Tue, 03 Nov 2026 12:12:21 GMT
Server: Kestrel
\end{lstlisting}

  \slidegap
  \aside{Сервис отвечает пустым 200 на всё. Один из самых частых багов первой
    лабораторной с middleware.}
\end{frame}

% 63
\begin{frame}{Сколько стоит middleware}
  \slidetop
  \claim{Цепочка дешёвая. Дорого то, что \hi{внутри}.}

  \begin{cbul}{accent}
    \item выполняется на \hi{каждом} запросе без исключения
    \item синхронное ожидание --- поток занят впустую
    \item тело --- поток, читается \lo{один раз}: прочитали для лога ---
      обработчику не осталось
    \item буферизация разрешает перемотку, но стоит памяти на каждый запрос
  \end{cbul}

  \slidegap
  \aside{Пишите middleware так, чтобы он стоил микросекунды.}
\end{frame}

% 64
\begin{frame}[plain]
  \keyline{Конвейер --- цепочка делегатов, и порядок --- это семантика.\\
    Ошибки ловятся в одном месте, снаружи всех.}
\end{frame}

% ======================================================== 4. HTTPCONTEXT ==
\section{4. HttpContext}

% 65
\begin{frame}[plain]
  \sectionopener{4}{HttpContext}{Всё о запросе в \hi{одном} объекте.}
\end{frame}

% 66
\begin{frame}{Один объект на запрос}
  \slidetop
  \begin{center}
    \begin{tikzpicture}[font=\footnotesize,
      pb/.style={box,minimum width=2.75cm,minimum height=0.75cm,inner sep=2pt,
                 font=\footnotesize\ttfamily}]
      \node[pb,draw=accent] (p1) at (0,0)     {Request};
      \node[pb,draw=accent] (p2) at (3.05,0)  {Response};
      \node[pb] (p3) at (6.1,0)   {User};
      \node[pb] (p4) at (9.15,0)  {RequestAborted};
      \node[pb] (p5) at (0,-1.0)    {Items};
      \node[pb] (p6) at (3.05,-1.0) {Features};
      \node[pb] (p7) at (6.1,-1.0)  {TraceIdentifier};
      \node[pb] (p8) at (9.15,-1.0) {RequestServices};
      \node[pb] (p9) at (4.575,-2.0) {Connection};
      \node[draw=muted,line width=0.9pt,rounded corners=3pt,inner sep=8pt,
            fit=(p1)(p4)(p9)] (hc) {};
      \node[anchor=south west,font=\scriptsize\bfseries\ttfamily,text=muted]
        at ([yshift=2pt]hc.north west) {HttpContext};
    \end{tikzpicture}
  \end{center}

  \slidegap
  \aside{Создан сервером, когда запрос разобран; протянут через весь конвейер;
    выброшен, когда ответ отправлен. Живёт \hi{ровно} столько, сколько запрос.}
\end{frame}

% 67
\begin{frame}{\texttt{Request} и \texttt{Response}}
  \slidetop
  \versus{Request}{%
    \begin{cbul}{accent}
      \item метод, путь, параметры
      \item заголовки --- словарь
      \item тело --- поток, читается \hi{один раз}
    \end{cbul}}{Response}{%
    \begin{cbul}{alt}
      \item статус, заголовки
      \item тело --- поток, в который пишут
      \item \texttt{HasStarted}: заголовки ушли
    \end{cbul}}

  \slidegap
  \hrulethin{0.42}
  \rulegap
  {\large Сырой запрос уже \hi{разобран}: вы читаете свойства, а не строки.}
\end{frame}

% 68
\begin{frame}{Кто, ушёл ли, какой номер}
  \slidetop
  {\renewcommand{\arraystretch}{1.6}
  \begin{tabular}{@{}L{3.8cm}L{9.0cm}@{}}
    \hi{\texttt{User}}            & кто пришёл; пустой, пока аутентификация не
      заполнила \\
    \hi{\texttt{RequestAborted}}  & токен отмены: клиент закрыл вкладку, оборвалась
      сеть \\
    \hi{\texttt{TraceIdentifier}} & уникален для запроса: в ответе об ошибке и в
      каждой строке лога \\
  \end{tabular}}

  \slidegap
  \aside{\texttt{RequestAborted} протягивают через весь стек до запроса в базу:
    клиент ушёл --- вся цепочка \hi{остановилась}, а не дорабатывает впустую.}
\end{frame}

% 69
\begin{frame}{\texttt{Items}}
  \slidetop
  \versus{сервис «на запрос»}{%
    \begin{cbul}{accent}
      \item свой маленький тип
      \item получают через конструктор
      \item типы проверяет \hi{компилятор}
    \end{cbul}}{Items}{%
    \begin{cbul}{alt}
      \item словарь из \texttt{object} в \texttt{object}
      \item ключ --- строка, договорённость
      \item ошибка --- во время работы
    \end{cbul}}

  \slidegap
  \hrulethin{0.42}
  \rulegap
  {\large Одно-два значения --- \texttt{Items}. Что-то \hi{структурное} --- свой
    тип.}
\end{frame}

% 70
\begin{frame}{Остальное}
  \slidetop
  {\renewcommand{\arraystretch}{1.6}
  \begin{tabular}{@{}L{3.8cm}L{9.0cm}@{}}
    \texttt{Features}               & низкий уровень сервера: через него работают
      \texttt{Request} и \texttt{Response} \\
    \hi{\texttt{RequestServices}}   & область контейнера на этот запрос:
      освобождается вместе с ним \\
    \texttt{Connection}             & адрес и порт того, кто открыл соединение \\
  \end{tabular}}

  \slidegap
  \aside{Адрес --- точка выхода в сеть, а \lo{не человек}: за одним адресом
    бывает тысяча пользователей.}
\end{frame}

% 71
\begin{frame}[plain]
  \vote{Сервису заказов нужен идентификатор текущего пользователя. Передать
    ему \texttt{HttpContext} или идентификатор?}
\end{frame}

% 72
\begin{frame}{Три цены}
  \slidetop
  \claim{Контекст в бизнес-логике прибивает её \lo{к вебу}.}

  {\renewcommand{\arraystretch}{1.6}
  \begin{tabular}{@{}L{4.4cm}L{8.4cm}@{}}
    \lo{\textbf{только внутри HTTP}} & ни из фоновой задачи, ни из теста, ни из
      утилиты \\
    \lo{\textbf{живёт с запросом}}   & изменяемый, не потокобезопасный; ссылка
      после ответа --- мусор \\
    \lo{\textbf{сигнатура врёт}}     & что метод берёт из контекста --- не видно
      без чтения тела \\
  \end{tabular}}
\end{frame}

% 73
\begin{frame}[fragile]{Контекст не уходит дальше эндпоинта}
  \slidetop
{\renewcommand{\arraystretch}{1.5}
  \begin{tabular}{@{}L{1.0cm}l@{}}
    \lo{\textbf{нет}} & {\codexs Task<bool> IsOwnerAsync(Order order, HttpContext context);} \\
    \hi{\textbf{да}}  & {\codexs Task<bool> IsOwnerAsync(Order order, string userId, CancellationToken ct);} \\
  \end{tabular}}

  \slidegap
  \begin{center}
    \begin{tikzpicture}[font=\footnotesize]
      \node[box,draw=accent,minimum width=3.4cm] (e) at (0,0)
        {эндпоинт\\{\scriptsize\ttfamily HttpContext}};
      \node[box,minimum width=3.4cm] (sv) at (5.0,0)
        {сервис\\{\scriptsize\ttfamily userId, ct}};
      \node[box,minimum width=3.4cm] (rp) at (10.0,0)
        {репозиторий\\{\scriptsize\ttfamily ct}};
      \draw[arra] (e) -- node[above,font=\scriptsize,text=accent] {значения} (sv);
      \draw[arr] (sv) -- (rp);
    \end{tikzpicture}
  \end{center}

  \slidegap
  \aside{Бизнес-логика получает строку, число, токен --- и \hi{не знает}, что
    где-то есть HTTP.}
\end{frame}

% 74
\begin{frame}[fragile]{«Текущий пользователь»}
  \vspace{0.3em}
\begin{lstlisting}[style=csharp,basicstyle=\codesm]
interface ICurrentUser { string Id { get; } }

sealed class HttpCurrentUser(IHttpContextAccessor accessor) : ICurrentUser
{
    public string Id =>
        accessor.HttpContext?.User.FindFirstValue(ClaimTypes.NameIdentifier)
        ?? throw new InvalidOperationException("No user");
}
\end{lstlisting}

  \slidegap
  {\renewcommand{\arraystretch}{1.3}\small
  \begin{tabular}{@{}L{3.4cm}L{9.4cm}@{}}
    в вебе     & читает контекст \\
    в тесте    & \hi{заглушка} с нужным пользователем \\
    в фоне     & системный пользователь \\
  \end{tabular}}

  \vspace{0.4em}
  \aside{\texttt{IHttpContextAccessor} --- для инфраструктуры. В бизнес-логике
    --- \lo{сигнал на ревью}.}
\end{frame}

% 75
\begin{frame}{Контекст и фоновая работа}
  \slidetop
  \versus{скопировать значения}{%
    \begin{cbul}{accent}
      \item \texttt{userId}, \texttt{traceId} --- до запуска
      \item живут сколько нужно
    \end{cbul}}{захватить context}{%
    \begin{cbul}{alt}
      \item ответ ушёл --- объект отдан другому запросу
      \item мусор или исключение
    \end{cbul}}

  \slidegap
  \hrulethin{0.42}
  \rulegap
  {\large Та же ловушка \hi{захвата}, что с лямбдами, --- только у объекта есть
    владелец.}
\end{frame}

% 76
\begin{frame}[plain]
  \keyline{HttpContext --- всё о запросе, и поэтому его не носят вглубь.\\
    Берите нужное на границе.}
\end{frame}

% ======================================================== 5. BACKGROUND WORK ==
\section{5. Фоновая работа}

% 77
\begin{frame}[plain]
  \sectionopener{5}{Фоновая работа в том же хосте}{Хост запускает \hi{не
    только} сервер.}
\end{frame}

% 78
\begin{frame}[fragile]{\texttt{IHostedService}}
  \slidetop
  \claim{Два метода: \hi{запустить} и остановить.}

\begin{lstlisting}[style=csharp,basicstyle=\ttfamily\small]
public interface IHostedService
{
    Task StartAsync(CancellationToken cancellationToken);
    Task StopAsync(CancellationToken cancellationToken);
}
\end{lstlisting}

  \slidegap
  \begin{cbul}{accent}
    \item старт --- в порядке регистрации, остановка --- в \hi{обратном}
    \item веб-сервер запускается так же: это тоже хостед-сервис
    \item очистка раз в час, отложенные письма, ночной пересчёт, очередь
  \end{cbul}
\end{frame}

% 79
\begin{frame}[fragile]{Подвох старта}
  \slidetop
  \claim{Хост ждёт \hi{каждый} \texttt{StartAsync}, по очереди.}

\begin{lstlisting}[style=term,basicstyle=\termsm]
12:13:24.849 info: SlowStart[0]
      StartAsync: loading for 5 s
12:13:29.856 info: SlowStart[0]
      StartAsync done
12:13:29.885 info: Microsoft.Hosting.Lifetime[14]
      Now listening on: http://localhost:5021
\end{lstlisting}

  \slidegap
  {\renewcommand{\arraystretch}{1.4}
  \begin{tabular}{@{}L{4.4cm}L{8.4cm}@{}}
    \lo{\texttt{StartAsync} грузит 5 с} & сервер не слушает 5 с \\
    \hi{\texttt{BackgroundService}}     & в .NET 10 весь \texttt{ExecuteAsync} идёт в
      фоне: запуск \hi{не ждёт} даже код до первого \texttt{await} \\
  \end{tabular}}

  \slidegap
  \aside{Долгая инициализация --- внутри фоновой работы, а не в запуске.}
\end{frame}

% 80
\begin{frame}[fragile]{\texttt{BackgroundService}}
  \vspace{0.3em}
\begin{lstlisting}[style=csharp,basicstyle=\ttfamily\fontsize{6.8}{7.8}\selectfont,escapechar=|]
sealed class CleanupService(ILogger<CleanupService> logger) : BackgroundService
{
    protected override async Task ExecuteAsync(CancellationToken stoppingToken)
    {
        using var timer = new PeriodicTimer(TimeSpan.FromMinutes(5));
|\mk{b1}|        while (await timer.WaitForNextTickAsync(stoppingToken))
        {
|\mk{b2}|            try
            {
                logger.LogInformation("Cleanup started");
                // work goes here: delete stale rows, recount, send
            }
|\mk{b3}|            catch (Exception e) when (e is not OperationCanceledException)
            {
                logger.LogError(e, "Cleanup iteration failed");
            }
        }
    }
}

|\mk{b4}|// builder.Services.AddHostedService<CleanupService>();
\end{lstlisting}
  \begin{tikzpicture}[remember picture,overlay,every node/.style={rawlab,font=\small\bfseries,text=accent}]
    \foreach \m/\n in {b1/1,b2/2,b3/3} \node at ([xshift=11.9cm]\m) {\n};
    \node[font=\footnotesize,text=muted] at ([xshift=8.0cm]b4) {регистрация одной строкой};
  \end{tikzpicture}
\end{frame}

% 81
\begin{frame}{Четыре решения}
  \slidetop
  \claim{Каждое --- чей-то \lo{инцидент}.}

  {\renewcommand{\arraystretch}{1.5}
  \begin{tabular}{@{}L{4.4cm}L{8.4cm}@{}}
    \hi{\textbf{1\enspace PeriodicTimer с токеном}} & поток не занят; на остановке ожидание
      прерывается \\
    \hi{\textbf{2\enspace try внутри цикла}}        & упала итерация --- следующая всё
      равно будет \\
    \hi{\textbf{3\enspace фильтр на отмену}}        & отмена --- сигнал, а не ошибка: не
      глотаем \\
    \lo{\textbf{4\enspace исключение наружу}}       & хост останавливает \lo{всё}
      приложение \\
  \end{tabular}}

  \slidegap
  \aside{\texttt{try} вокруг цикла --- первая же ошибка завершит сервис
    навсегда, и вы узнаете через месяц.}
\end{frame}

% 82
\begin{frame}[fragile]{Исключение из \texttt{ExecuteAsync}}
  \slidetop
  \claim{Одна ошибка в фоне --- и \lo{сервис лежит}.}

\begin{lstlisting}[style=term,basicstyle=\termsm]
fail: Microsoft.Extensions.Hosting.Internal.Host[9]
      BackgroundService failed
      System.InvalidOperationException: boom
crit: Microsoft.Extensions.Hosting.Internal.Host[10]
      The HostOptions.BackgroundServiceExceptionBehavior is configured to StopHost. ...
      System.InvalidOperationException: boom
info: Microsoft.Hosting.Lifetime[0]
      Application is shutting down...
\end{lstlisting}

  \slidegap
  \aside{Так сделано намеренно: упавший фон не должен оставаться незамеченным.
    Код выхода при этом --- \lo{0}. Поведение меняется настройкой хоста --- лучше
    не менять.}
\end{frame}

% 83
\begin{frame}[fragile]{Зависимости «на запрос» в фоне}
  \vspace{0.3em}
\begin{lstlisting}[style=csharp,basicstyle=\ttfamily\fontsize{6.8}{7.8}\selectfont,escapechar=|]
sealed class ReportService(IServiceScopeFactory scopes) : BackgroundService
{
    protected override async Task ExecuteAsync(CancellationToken stoppingToken)
    {
        using var timer = new PeriodicTimer(TimeSpan.FromMinutes(5));

        while (await timer.WaitForNextTickAsync(stoppingToken))
        {
|\mk{r1}|            using var scope = scopes.CreateScope();
|\mk{r2}|            var reports = scope.ServiceProvider.GetRequiredService<ReportBuilder>();
            await reports.BuildAsync(stoppingToken);
        }
    }
}
\end{lstlisting}

  \rulegap
  {\small В конструкторе, в Development: \texttt{\lo{Cannot consume scoped service
    'ReportBuilder' from singleton}}\par}
  \rulegap
  \aside{Одна итерация --- \hi{одна область}, один контекст базы.}
\end{frame}

% 84
\begin{frame}[fragile]{Расписание: Quartz.NET}
  \slidetop
  \begin{center}
    \begin{tikzpicture}[font=\footnotesize,
      qb/.style={box,minimum width=3.6cm,minimum height=1.2cm}]
      \node[qb] (j) at (0,0)   {\textbf{Job}\\что делать};
      \node[qb] (t) at (4.3,0) {\textbf{Trigger}\\когда};
      \node[qb,draw=accent,line width=1.0pt] (s) at (8.6,0)
        {\textbf{Scheduler}\\кто следит};
    \end{tikzpicture}
  \end{center}

  \slidegap
  \begin{center}
  {\ttfamily\fontsize{14}{16}\selectfont
  \begin{tabular}{cccccc}
    0 & 0 & 3 & ? & * & \hi{MON-FRI} \\[-0.2em]
    {\scriptsize\rmfamily\color{muted} сек} & {\scriptsize\rmfamily\color{muted} мин}
    & {\scriptsize\rmfamily\color{muted} час} & {\scriptsize\rmfamily\color{muted} день}
    & {\scriptsize\rmfamily\color{muted} месяц} & {\scriptsize\rmfamily\color{muted}
      дни недели} \\
  \end{tabular}}
  \end{center}

  \slidegap
  \aside{В три ночи по будням. Секунды --- \hi{первым} полем, в отличие от
    классического cron. Scheduler сам --- хостед-сервис.}
\end{frame}

% 85
\begin{frame}{Чего у таймера не будет}
  \slidetop
  {\renewcommand{\arraystretch}{1.6}
  \begin{tabular}{@{}L{4.2cm}L{8.6cm}@{}}
    \hi{\textbf{пропущенные запуски}} & остановлено в три ночи --- в девять утра
      решить: запустить или забыть \\
    \hi{\textbf{данные задания}}      & параметры, которые триггер передаёт при
      запуске \\
    \hi{\textbf{пауза}}               & приостановить и возобновить триггер на
      ходу \\
    \hi{\textbf{хранение в базе}}     & задания и триггеры переживают перезапуск \\
  \end{tabular}}

  \slidegap
  \aside{Писать это на таймере --- значит писать \lo{свой} планировщик.}
\end{frame}

% 86
\begin{frame}[fragile]{Задание и триггер}
  \vspace{0.3em}
\begin{lstlisting}[style=csharp,basicstyle=\ttfamily\fontsize{6.8}{7.7}\selectfont,escapechar=|]
|\mk{q1}|[DisallowConcurrentExecution]
sealed class NightlyReportJob : IJob
{
    public ValueTask Execute(IJobExecutionContext context,
        CancellationToken cancellationToken)
    {
        // build and send the report; cancellationToken is the stop signal
        return ValueTask.CompletedTask;
    }
}

builder.Services.AddQuartz(q =>
{
|\mk{q2}|    var key = new JobKey("nightly-report");
    q.AddJob<NightlyReportJob>(j => j.WithIdentity(key));
|\mk{q3}|    q.AddTrigger(t => t.ForJob(key).WithCronSchedule("0 0 3 ? * MON-FRI"));
});
|\mk{q4}|builder.Services.AddQuartzHostedService(o => o.WaitForJobsToComplete = true);
\end{lstlisting}
  \begin{tikzpicture}[remember picture,overlay,every node/.style={rawlab,font=\small\bfseries,text=accent}]
    \foreach \m/\n in {q1/1,q2/2,q3/3,q4/4} \node at ([xshift=11.9cm]\m) {\n};
  \end{tikzpicture}

  \vspace{0.1em}
  \aside{\hi{1} не параллельно \quad \hi{2} ключ задания \quad \hi{3} cron
    \quad \hi{4} дождаться заданий на остановке}
\end{frame}

% 87
\begin{frame}[fragile]{Три ночи --- по какому времени?}
  \slidetop
  \claim{В контейнере --- \lo{UTC}: 03:00 UTC --- это 06:00 в Минске.}

\begin{lstlisting}[style=csharp,basicstyle=\codexs]
q.AddTrigger(t => t.ForJob(key)
    .WithCronSchedule("0 0 3 ? * MON-FRI", x => x.InTimeZone(
        TimeZoneInfo.FindSystemTimeZoneById("Europe/Minsk"))));
\end{lstlisting}

  \slidegap
  \begin{cbul}{accent}
    \item перевод часов: час повторяется или пропадает --- запуск \lo{дважды}
      или ни разу
    \item расписание для людей --- с \hi{явным} часовым поясом
    \item остальное --- проще думать в UTC
  \end{cbul}
\end{frame}

% 88
\begin{frame}{Наложение запусков}
  \slidetop
  \claim{Триггер каждые 30 минут, работа --- \lo{40}.}

  \begin{center}
    \begin{tikzpicture}[font=\scriptsize,x=0.075cm,
      rb/.style={box,minimum height=0.5cm,inner sep=1pt,anchor=west}]
      \draw[hairline,line width=0.6pt] (0,0) -- (120,0);
      \foreach \x in {0,30,60,90,120} {
        \draw[hairline] (\x,0.08) -- (\x,-0.08);
        \node[text=muted,below=2pt] at (\x,0) {\x};
      }
      \node[text=muted,anchor=west] at (122,0) {мин};
      \node[anchor=east,text=muted] at (-3,2.1)  {запуск 1};
      \node[anchor=east,text=alt]   at (-3,1.35) {без атрибута};
      \node[anchor=east,text=accent] at (-3,0.6) {\texttt{[DisallowConcurrentExecution]}};
      \node[rb,minimum width=3.0cm] at (0,2.1) {40 мин};
      \node[rb,draw=alt,line width=1.0pt,minimum width=3.0cm] at (30,1.35)
        {\textcolor{alt}{запуск 2}};
      \node[rb,draw=accent,line width=1.0pt,minimum width=3.0cm] at (40,0.6)
        {\textcolor{accent}{запуск 2}};
      \fill[alt,opacity=0.15] (30,1.05) rectangle (40,2.4);
      \node[text=alt,anchor=west] at (72,1.35) {два одновременно, одни данные};
      \node[text=accent,anchor=west] at (82,0.6) {ждёт первый};
    \end{tikzpicture}
  \end{center}

  \slidegap
  \aside{У \texttt{BackgroundService} с \texttt{await} в цикле наложения нет:
    следующая итерация не начнётся до конца текущей.}
\end{frame}

% 89
\begin{frame}[plain]
  \vote{Два экземпляра сервиса, в каждом задание «в три ночи отправить счета».
    Сколько раз клиенты получат счета?}
\end{frame}

% 90
\begin{frame}{Три пути к повтору}
  \slidetop
  \claim{Работа запустится \lo{больше} одного раза. Вопрос --- когда.}

  {\renewcommand{\arraystretch}{1.6}
  \begin{tabular}{@{}L{4.6cm}L{8.2cm}@{}}
    \lo{\textbf{два экземпляра}}       & у каждого свой планировщик --- оба честно
      сработают \\
    \lo{\textbf{остановка посередине}} & таймаут вышел, процесс убит, после
      перезапуска --- снова \\
    \lo{\textbf{наложение}}            & предыдущий запуск ещё идёт, а триггер уже
      сработал \\
  \end{tabular}}

  \slidegap
  \aside{Клиенты получат по два счёта. Это не гипотетика --- это классика.}
\end{frame}

% 91
\begin{frame}{Идемпотентность}
  \slidetop
  \claim{Повторный запуск \hi{не делает} ничего лишнего.}

  {\renewcommand{\arraystretch}{1.6}
  \begin{tabular}{@{}L{3.2cm}L{9.6cm}@{}}
    \hi{\textbf{счёт}}      & помечается отправленным до или атомарно с отправкой
      --- повтор видит пометку \\
    \hi{\textbf{отчёт}}     & пишется по ключу «дата + тип» --- повтор
      перезаписывает тот же \\
    \hi{\textbf{сообщение}} & подтверждается после обработки --- повторная доставка
      безвредна \\
  \end{tabular}}

  \slidegap
  \aside{То же слово, что у методов HTTP, и та же причина: мир ненадёжен,
    повторы случаются.}
\end{frame}

% 92
\begin{frame}{Очередь}
  \slidetop
  \claim{Доставка --- «\hi{хотя бы} один раз».}

  \begin{center}
    \begin{tikzpicture}[font=\footnotesize,
      ob/.style={box,minimum width=3.0cm,minimum height=1.0cm}]
      \node[ob] (a) at (0,0)   {сервис};
      \node[ob] (q) at (5.0,0) {очередь};
      \node[ob,draw=accent,line width=1.0pt] (h) at (10.0,0) {обработчик};
      \draw[arr] (a) -- node[above=2pt,font=\scriptsize] {«отправить счёт»} (q);
      \draw[arra] ([yshift=0.15cm]q.east) -- node[above=2pt,font=\scriptsize,text=accent]
        {доставка} ([yshift=0.15cm]h.west);
      \draw[arrl] ([yshift=-0.15cm]q.east) -- node[below=2pt,font=\scriptsize,text=alt]
        {повтор} ([yshift=-0.15cm]h.west);
    \end{tikzpicture}
  \end{center}

  \slidegap
  \aside{«Ровно один раз» в распределённой системе стоит огромных денег. Значит,
    обработчик обязан быть \hi{идемпотентным} --- теперь это записано в контракте
    очереди.}
\end{frame}

% 93
\begin{frame}{Один запуск на всех}
  \slidetop
  {\renewcommand{\arraystretch}{1.6}
  \begin{tabular}{@{}L{4.6cm}L{8.2cm}@{}}
    \hi{\textbf{кластерный режим}}         & Quartz хранит задания в общей базе,
      экземпляры договариваются \\
    \hi{\textbf{распределённая блокировка}} & перед стартом захватить запись с
      коротким сроком жизни \\
  \end{tabular}}

  \slidegap
  \hrulethin{0.42}
  \rulegap
  {\large Ни то ни другое \lo{не отменяет} идемпотентности.}

  \vspace{0.6em}
  \aside{Захвативший может умереть посередине --- остановка никуда не делась.}
\end{frame}

% 94
\begin{frame}{Где жить фоновой работе}
  \slidetop
  \versus{тот же процесс}{%
    \begin{cbul}{accent}
      \item лёгкая и редкая
      \item общие конфигурация, контейнер, логгер
    \end{cbul}}{отдельный рабочий сервис}{%
    \begin{cbul}{alt}
      \item тяжёлая, занимает ядра на час
      \item тот же код, другой хост --- без сервера
    \end{cbul}}

  \slidegap
  \hrulethin{0.42}
  \rulegap
  {\large Пора выносить --- задержки запросов \hi{повторяют} график фоновой
    работы.}

  \vspace{0.6em}
  \aside{Фон делит с сервером пул потоков: тяжёлый пересчёт замедляет ответы.}
\end{frame}

% 95
\begin{frame}{Фоновая работа невидима}
  \slidetop
  \claim{Упавший запрос видит клиент. Переставшее запускаться задание ---
    \lo{никто}.}

  \begin{cbul}{accent}
    \item лог \hi{начала и конца} с длительностью
    \item метрика «последний успешный запуск»
    \item оповещение: отчёт не строился двое суток --- кто-то узнаёт утром
  \end{cbul}
\end{frame}

% 96
\begin{frame}[plain]
  \keyline{Фоновая работа запускается столько раз, сколько экземпляров и
    перезапусков.\\ Пишите её так, чтобы повтор был безвреден.}
\end{frame}

% 97
\begin{frame}{Итог}
  \slidetop
  \begin{center}
    \begin{tikzpicture}[font=\footnotesize,
      hb/.style={box,minimum width=2.3cm,minimum height=0.9cm}]
      \node[hb] (c1) at (0,0)    {конфигурация};
      \node[hb] (c2) at (2.6,0)  {контейнер};
      \node[hb] (c3) at (5.2,0)  {логгер};
      \node[hb,draw=accent,line width=1.0pt] (c4) at (7.8,0) {сервер};
      \node[hb] (c5) at (10.4,0) {фоновые\\сервисы};
      \node[draw=muted,line width=0.9pt,rounded corners=3pt,inner sep=8pt,
            fit=(c1)(c5)] (h) {};
      \node[anchor=south west,font=\scriptsize\bfseries,text=muted]
        at ([yshift=2pt]h.north west) {ХОСТ};
      \foreach \i/\t in {1/ошибки,2/логирование,3/маршрутизация,4/авторизация} {
        \node[box,minimum width=2.3cm,minimum height=0.7cm,inner sep=2pt]
          (k\i) at ({(\i-1)*2.6},-2.1) {\t};
      }
      \node[box,draw=accent,line width=1.0pt,minimum width=2.3cm,minimum height=0.7cm,
            inner sep=2pt] (k5) at (10.4,-2.1) {обработчик};
      \foreach \i [evaluate=\i as \j using int(\i+1)] in {1,...,4}
        \draw[arra] (k\i) -- (k\j);
      \node[anchor=north west,font=\scriptsize\bfseries,text=muted]
        at ([yshift=-4pt]k1.south west) {КОНВЕЙЕР: НА КАЖДЫЙ ЗАПРОС};
    \end{tikzpicture}
  \end{center}

  \slidegap
  \aside{Дальше --- то, что стоит в конце конвейера: обработчики, параметры
    из запроса, ответ в JSON.}
\end{frame}

% 98
\begin{frame}{Лабораторная 8}
  \slidetop
  \claim{Первый сервис --- он живёт \hi{до конца} курса.}

  \begin{minipage}[t]{0.50\textwidth}
    \begin{cbul}{accent}
      \item шаблон \texttt{web}, минимальный \texttt{Program.cs}
      \item middleware логирования --- класс
      \item middleware ошибок --- класс, \hi{первый}
      \item \texttt{/health} и бросающий эндпоинт
      \item \texttt{BackgroundService} с \texttt{PeriodicTimer}
    \end{cbul}
  \end{minipage}\hfill
  \begin{minipage}[t]{0.44\textwidth}
    \eyebrow{проверка}\par\vspace{0.5em}
    \begin{cbul}{alt}
      \item \texttt{curl}: здоровье, обычный, бросающий
      \item в логе строка на запрос и трассировка на ошибку
      \item Ctrl+C --- остановка за секунды
    \end{cbul}
  \end{minipage}
\end{frame}

\end{document}
